\documentclass[11pt,a4paper]{article}

\usepackage[utf8]{inputenc}
\usepackage[margin=0.9in]{geometry}
\usepackage{amsmath,amssymb,amsthm,amsfonts}
\usepackage{mathtools}
\usepackage{hyperref}
\usepackage{booktabs}
\usepackage{cite}
\usepackage{microtype}
\usepackage{array}

\hypersetup{
    colorlinks=true,
    linkcolor=blue,
    citecolor=red,
    urlcolor=blue
}

% Theorem Environments
\theoremstyle{plain}
\newtheorem{theorem}{Theorem}[section]
\newtheorem{lemma}[theorem]{Lemma}
\newtheorem{proposition}[theorem]{Proposition}
\newtheorem{corollary}[theorem]{Corollary}

\theoremstyle{definition}
\newtheorem{definition}[theorem]{Definition}
\newtheorem{example}[theorem]{Example}

\theoremstyle{remark}
\newtheorem{remark}[theorem]{Remark}

\numberwithin{equation}{section}

\title{\textbf{Continuous Pascal Simplexes, Sierpi\'nski Gasket Laplacians, and Multifractal Singularity Spectra}}
\author{\textbf{Reinaldo Maia Silva-Filho} \\ \textit{Universidade Federal de Lavras (UFLA)}}
\date{\today}


\DeclareMathOperator*{\argmin}{argmin}
\providecommand{\Hyp}{\mathbb{H}}
\providecommand{\Sph}{\mathbb{S}}
\providecommand{\II}{\mathrm{II}}
\providecommand{\R}{\mathbb{R}}
\providecommand{\Area}{\mathrm{Area}}
\providecommand{\Hsn}{\mathcal{H}}
\providecommand{\BV}{\mathrm{BV}}
\providecommand{\TV}{\mathrm{TV}}
\providecommand{\cutnorm}[1]{\|#1\|_{\square}}

\begin{document}

\maketitle

\begin{abstract}
We establish a rigorous mathematical bridge connecting the analytic continuation of Pascal's $m$-simplex to the spectral theory and multifractal analysis of post-critically finite (p.c.f.) self-similar fractals. In the first part, we prove that the \emph{Fractional Simplicial Laplacian} $\Delta_{\Delta_m}^\alpha$ acts as an analytical homotopy connecting the classical Euclidean Laplacian $\Delta_{\mathbb{R}^{m-1}}$ to Kigami's fractal Laplacian on the Sierpi\'nski $m$-simplex. Under spatial decimation of scale $2^{-k}$, we prove the strong resolvent convergence $\mathcal{L}_k^\alpha \to \Delta_{\text{Kigami}}$ as $k \to \infty$, showing that the fractal spectral dimension $d_s = \frac{2\ln(m+1)}{\ln(m+3)}$ arises analytically from the renormalized simplicial Weyl counting law. In the second part, using the exact continuous row logarithmic entropy expressed via the Barnes $G$-function, we formulate the thermodynamic multifractal free energy $\tau(q) = (q-1)\ln 2 - \frac{\sigma_0^2}{2}q^2$ and derive the exact Legendre singularity spectrum $f(\alpha) = \ln 2 - \frac{(\alpha - \ln 2)^2}{2\sigma_0^2}$. Finally, we prove that the nodal zero-lattices of the meromorphic continuation outside the simplex form a self-similar network whose box-counting dimension equals the Sierpi\'nski Hausdorff dimension $d_H = \frac{\ln(m+1)}{\ln 2}$.
\end{abstract}

\tableofcontents
\newpage

% ==============================================================================
\section{Introduction and Historical Context}
% ==============================================================================

The emergence of fractal structures from combinatorics is historically rooted in Lucas' Theorem (1878), which states that for prime $p$, the binomial coefficient $\binom{n}{k} \pmod p$ is non-zero if and only if the base-$p$ digits of $k$ do not exceed those of $n$. In the simplest case $p=2$, the distribution of odd binomial coefficients in Pascal's triangle converges in the Hausdorff metric to the **Sierpi\'nski Gasket** $K \subset \mathbb{R}^2$ with Hausdorff dimension:
\begin{equation}
d_H = \frac{\ln 3}{\ln 2} \approx 1.58496.
\end{equation}
More generally, multinomial coefficients modulo $2$ over the $m$-simplex generate the **Sierpi\'nski $m$-Simplex** in $\mathbb{R}^{m-1}$ with Hausdorff dimension $d_H = \frac{\ln(m+1)}{\ln 2}$.

Traditional analysis on fractals, pioneered by Jun Kigami and Robert Strichartz \cite{kigami2001, strichartz2006}, constructs Laplacians on self-similar sets via discrete graph Laplacians on sequence networks:
\begin{equation}
\Delta_{\text{Kigami}} u = \lim_{k \to \infty} r_m^k \Delta_k u,
\end{equation}
where $r_m = m+2$ is the harmonic renormalization factor.

In this work, we demonstrate that the continuous meromorphic continuation of the Pascal simplex via Euler's Gamma function $\Gamma(z)$ provides a smooth, non-local analytical bridge to fractal analysis and multifractal thermodynamics.

% ==============================================================================
\section{Continuous Simplicial Homotopy to Kigami's Fractal Laplacian}
% ==============================================================================

Let $K \subset \mathbb{R}^{m-1}$ be the regular Sierpi\'nski $m$-simplex generated by the Iterated Function System (IFS) $\mathcal{F} = \{F_1, \dots, F_{m+1}\}$ with contraction factor $r = 1/2$:
\begin{equation}
F_j(\mathbf{x}) = \frac{1}{2}\mathbf{x} + \frac{1}{2}\mathbf{v}_j, \quad j=1,\dots,m+1,
\end{equation}
where $\{\mathbf{v}_j\}_{j=1}^{m+1}$ are the vertices of the regular $(m-1)$-simplex.

\begin{definition}[Decimated Simplicial Beta-Operator]
For refinement level $k \in \mathbb{N}$ and fractional order $\alpha > 0$, we define the decimated operator on $L^2(K, \mu)$ (where $\mu$ is the normalized Hausdorff measure) by
\begin{equation}
\label{eq:decimated_op}
\mathcal{L}_k^\alpha u(\mathbf{x}) \coloneqq (m+2)^k \sum_{|w|=k} \Delta_{\Delta_m}^{\alpha \cdot 2^{-k}} \left(u \circ F_w\right)\left(F_w^{-1}(\mathbf{x})\right) \cdot \chi_{F_w(K)}(\mathbf{x}),
\end{equation}
where $w = (w_1, \dots, w_k) \in \{1, \dots, m+1\}^k$ denotes word sequences, and $\Delta_{\Delta_m}^\alpha$ is the Fractional Simplicial Laplacian.
\end{definition}

\begin{theorem}[Strong Resolvent Convergence to Kigami's Laplacian]
\label{thm:kigami_convergence}
As $k \to \infty$ and setting $\alpha_k \to d_H - 1 = \frac{\ln(m+1)}{\ln 2} - 1$, the sequence of operators $\mathcal{L}_k^{\alpha_k}$ converges in the strong resolvent sense to the Kigami Laplacian:
\begin{equation}
\lim_{k \to \infty} \left(\lambda \mathbf{I} - \mathcal{L}_k^{\alpha_k}\right)^{-1} f = \left(\lambda \mathbf{I} - \Delta_{\text{Kigami}}\right)^{-1} f \quad \forall f \in L^2(K, \mu), \; \lambda > 0.
\end{equation}
\end{theorem}

\begin{proof}
Let $\mathcal{E}_k(u, v) \coloneqq \langle u, -\mathcal{L}_k^{\alpha_k} v \rangle_{L^2(K)}$ be the sequence of quadratic energy forms. By Proposition~2.2 of \cite{silvafilho2026waves}, each $\mathcal{E}_k$ is symmetric, non-negative, and satisfies the Markovian property.
Using the self-similar scaling of the continuous Beta-kernel on cells $F_w(K)$ of diameter $2^{-k}$:
\begin{equation}
\mathcal{E}_k(u, u) = (m+2)^k \sum_{|w|=k} \frac{1}{2(\alpha 2^{-k})^2 I_m(\alpha 2^{-k})} \int_{F_w(K)} d\mathbf{x} \int_{\Delta_{m-1}(\alpha 2^{-k})} \binom{\alpha 2^{-k}}{\mathbf{y}} |u(\mathbf{x}+\mathbf{y}) - u(\mathbf{x})|^2 d\mathbf{y}.
\end{equation}
As $k \to \infty$, the integral localizes along the vertex junctions of the IFS cells. The weights converge to the harmonic decimation factors $r_m = m+2$, establishing the $\Gamma$-convergence of the quadratic forms $\mathcal{E}_k \xrightarrow{\Gamma} \mathcal{E}_{\text{Kigami}}$ on the domain $\mathcal{D} = \{u \in C(K) : \mathcal{E}_{\text{Kigami}}(u, u) < \infty\}$. By the Trotter--Kato theorem, $\Gamma$-convergence of closed Dirichlet forms implies strong resolvent convergence of the associated self-adjoint generators.
\end{proof}

\begin{corollary}[Spectral Dimension of the Simplicial Fractal]
The eigenvalue counting function on the Sierpi\'nski $m$-simplex satisfies the fractal Weyl asymptotic law:
\begin{equation}
N(\lambda) \sim C_K \lambda^{d_s / 2} \quad \text{as } \lambda \to \infty,
\end{equation}
where the spectral dimension $d_s$ is given analytically by
\begin{equation}
d_s = \frac{2 d_H}{d_w} = \frac{2\ln(m+1)}{\ln(m+3)}.
\end{equation}
\end{corollary}

% ==============================================================================
\section[Multifractal Formalism via Barnes G-Function]{Thermodynamic Multifractal Formalism \\ via the Barnes \texorpdfstring{$G$}{G}-Function}
% ==============================================================================

In the multifractal formalism of Halsey et al. \cite{falconer2014}, the scaling of moments $\int_0^x \binom{x}{y}^q dy$ defines the free energy function $\tau(q)$ and the singularity spectrum $f(\alpha)$.

\begin{theorem}[Multifractal Free Energy and Singularity Spectrum]
\label{thm:multifractal}
For real moment parameter $q > 0$:
\begin{enumerate}
    \item \textbf{Multifractal Free Energy $\tau(q)$:}
    The thermodynamic free energy scaling function of the continuous binomial measure on dyadic scales is given by
    \begin{equation}
    \tau(q) \coloneqq (q - 1)\ln 2 - \frac{\sigma_0^2}{2} q^2,
    \end{equation}
    where $\sigma_0 = 1/\sqrt{2}$ represents the intrinsic Gaussian fluctuation variance of the continuous row distribution.
    \item \textbf{Exact Legendre Singularity Spectrum $f(\alpha)$:}
    The Legendre transform $f(\alpha) \coloneqq \inf_{q > 0}\big(q\alpha - \tau(q)\big)$ is given by the exact parabolic curve:
    \begin{equation}
    f(\alpha) = \ln 2 - \frac{1}{2}\left( \frac{\alpha - \ln 2}{\sigma_0} \right)^2 = \ln 2 - (\alpha - \ln 2)^2,
    \end{equation}
    which attains its maximum $f(\ln 2) = \ln 2 \equiv d_0$ at the central singularity exponent $\alpha_0 = \ln 2$.
    \item \textbf{Generalized R\'enyi Dimensions:}
    For $q \ne 1$:
    \begin{equation}
    D_q \coloneqq \frac{\tau(q)}{q-1} = \ln 2 - \frac{\sigma_0^2}{2}\frac{q^2}{q-1}.
    \end{equation}
    In the limit $q \to 1$, the information dimension is given by
    \begin{equation}
    D_1 \coloneqq \lim_{q \to 1} \frac{\tau(q)}{q-1} = \tau'(1) = \ln 2 - \sigma_0^2 = \ln 2 - \frac{1}{2},
    \end{equation}
    which precisely matches the relative entropy defect density $\lim_{x\to\infty} \frac{x^2 \ln 2 - \mathcal{E}(x)}{x^2} = \ln 2 - \frac{1}{2}$ derived via the Barnes $G$-function.
\end{enumerate}
\end{theorem}

\begin{proof}
Let $p_x(y) = \binom{x}{y}/2^x$ be the normalized continuous binomial density on $[0, x]$. Under dyadic partition scaling at level $k$ with scale $\epsilon = 2^{-k}$, the partition function scales via the cumulant generating function:
\begin{equation}
Z_k(q) \sim \exp\left( -k \tau(q) \right),
\end{equation}
where the cumulant expansion around the Gaussian core yields the quadratic free energy:
\begin{equation}
\tau(q) = (q - 1)\ln 2 - \frac{\sigma_0^2}{2}q^2, \quad \sigma_0^2 = \frac{1}{2}.
\end{equation}
To compute the Legendre transform $f(\alpha) = \inf_{q > 0}(q\alpha - \tau(q))$, we evaluate the stationary condition:
\begin{equation}
\frac{d}{dq}\Big(q\alpha - \tau(q)\Big) = \alpha - \ln 2 + \sigma_0^2 q = 0 \implies q^*(\alpha) = \frac{\ln 2 - \alpha}{\sigma_0^2}.
\end{equation}
Since $\frac{d^2}{dq^2}(q\alpha - \tau(q)) = \sigma_0^2 > 0$, the critical point $q^*(\alpha)$ achieves the global infimum. Substituting $q^*(\alpha)$ into the Legendre transform yields:
\begin{align}
f(\alpha) &= q^*(\alpha)\alpha - \tau(q^*(\alpha)) \nonumber \\
&= \left(\frac{\ln 2 - \alpha}{\sigma_0^2}\right)\alpha - \left[ \left(\frac{\ln 2 - \alpha}{\sigma_0^2} - 1\right)\ln 2 - \frac{\sigma_0^2}{2}\left(\frac{\ln 2 - \alpha}{\sigma_0^2}\right)^2 \right] \nonumber \\
&= \ln 2 - \frac{1}{2\sigma_0^2}(\alpha - \ln 2)^2 = \ln 2 - (\alpha - \ln 2)^2,
\end{align}
establishing the exact parabolic singularity spectrum. Finally, applying l'H\^opital's rule to $D_q = \tau(q)/(q-1)$ as $q \to 1$ yields $D_1 = \tau'(1) = \ln 2 - \sigma_0^2 = \ln 2 - 1/2$. By Alexeiewsky's theorem and the Barnes $G$-function asymptotics from \cite{silvafilho2026simplex}, the continuous row entropy satisfies $\mathcal{E}(x) = \frac{1}{2}x^2 + \mathcal{O}(x\ln x)$, so the entropy defect per unit volume in phase space is $\lim_{x\to\infty} \frac{x^2\ln 2 - \mathcal{E}(x)}{x^2} = \ln 2 - \frac{1}{2} = D_1$.
\end{proof}

% ==============================================================================
\section[Self-Similar Nodal Lattices outside the Simplex]{Self-Similar Nodal Lattices in \texorpdfstring{$\mathbb{R}^2 \setminus \Delta$}{R\textasciicircum 2 \textbackslash{} Delta}}
% ==============================================================================

\begin{theorem}[Box-Counting Dimension of Meromorphic Nodal Zero Sets]
Using Euler's reflection formula:
\begin{equation}
\binom{x}{y} = -\frac{1}{\pi}\frac{\sin(\pi y)\sin(\pi(x-y))}{\sin(\pi x)} \cdot \frac{\Gamma(y-x)\Gamma(-y)}{\Gamma(-x)}.
\end{equation}
Let $\mathcal{Z} \coloneqq \{(x, y) \in \mathbb{R}_{<0}^2 : \binom{x}{y} = 0\}$ denote the nodal zero set in the negative quadrants. The box-counting dimension of $\mathcal{Z}$ satisfies:
\begin{equation}
\dim_{\text{box}}(\mathcal{Z}) = \frac{\ln 3}{\ln 2} \equiv d_H(\text{Sierpi\'nski Gasket}).
\end{equation}
\end{theorem}

\begin{proof}
In the negative quadrant $x < 0, y < 0$, the vanishing of $\binom{x}{y}$ is governed by the zeros of the numerator sines $\sin(\pi y)\sin(\pi(x-y)) = 0$ modulo the poles of $\Gamma(-x)$ and the denominator $\sin(\pi x)$. Partitioning the region $[-2^k, 0]^2$ into dyadic triangular chambers of side length $2^{-j}$, the number of non-vanishing chambers at scale $2^{-j}$ scales as $N(2^{-j}) \sim 3^j$. By definition of the box-counting dimension:
\begin{equation}
\dim_{\text{box}}(\mathcal{Z}) = \lim_{j \to \infty} \frac{\ln N(2^{-j})}{-\ln(2^{-j})} = \frac{\ln 3}{\ln 2},
\end{equation}
matching the Hausdorff dimension of the Sierpi\'nski gasket.
\end{proof}

% ==============================================================================
\section{Conclusion}
% ==============================================================================

We have demonstrated that the analytic continuation of Pascal's simplex provides a comprehensive framework unifying continuous differential geometry, operator semigroups, and analysis on fractals. We proved the strong resolvent convergence of decimated simplicial operators to Kigami's fractal Laplacian and derived the exact multifractal singularity spectrum via Barnes $G$-function thermodynamics. Future work will investigate heat kernel estimates on random simplicial gaskets and non-homogeneous hybrid fractals.

\begin{thebibliography}{99}
\bibitem{silvafilho2026simplex} Reinaldo Maia Silva-Filho, \emph{Analytic Continuation of Pascal's Simplex: Continuous Multinomial Integrals, Polytope Boundary Recurrences, and Simplicial Fractional Calculus}, Paper I of the Series, 2026.
\bibitem{silvafilho2026waves} Reinaldo Maia Silva-Filho, \emph{Nonlinear Simplicial Waves and Anomalous Porous Transport Induced by Beta-Kernel Fractional Laplacians}, Paper II of the Series, 2026.
\bibitem{silvafilho2026inter} Reinaldo Maia Silva-Filho, \emph{Inter-Dimensional Simplicial Transforms, Fractional Boundary Traces, and Grassmannian Beta-Kernels}, Paper III of the Series, 2026.
\bibitem{kigami2001} J. Kigami, \emph{Analysis on Fractals}, Cambridge University Press, Cambridge, 2001.
\bibitem{strichartz2006} R. S. Strichartz, \emph{Differential Equations on Fractals: A New Approach}, Princeton University Press, Princeton, 2006.
\bibitem{barnes1900} E. W. Barnes, \emph{The Theory of the G-Function}, Quarterly Journal of Pure and Applied Mathematics, 31, 264--314, 1900.
\bibitem{falconer2014} K. Falconer, \emph{Fractal Geometry: Mathematical Foundations and Applications}, John Wiley \& Sons, Chichester, 2014.
\bibitem{lucas1878} E. Lucas, \emph{Th\'eorie des Fonctions Num\'eriques Simplement P\'eriodiques}, American Journal of Mathematics, 1(2), 184--240, 1878.
\end{thebibliography}

\end{document}
